\chapter{Eksperymenty i wyniki}
\label{chap:results}
\section{Warunki wykonania eksperymentów}
\writingnote{Wpisz rzeczywiste konfiguracje, wersje, okresy danych, liczbę przebiegów i charakterystykę sprzętu. Wyniki z laptopa i Raspberry Pi raportuj oddzielnie.}
\section{Skuteczność metod i fałszywe alarmy}
\IfFileExists{generated/model-comparison.tex}{%
  \input{generated/model-comparison}
}{%
  \statusnote{Brak zatwierdzonych wyników. Tabela pojawi się po wyeksportowaniu rzeczywistych pomiarów do generated/model-comparison.tex.}
}
\writingnote{Przedstaw rozrzut wyników pomiędzy przebiegami i analizę błędów. Nie zastępuj brakujących pomiarów przykładowymi liczbami.}
\section{Transfer na dane fizyczne}
\writingnote{Omów wyniki testów fizycznego prototypu oraz różnice wobec danych treningowych.}
\section{Wpływ liczby węzłów}
\writingnote{Porównaj skuteczność i obciążenie przy różnych topologiach oraz natężeniu telemetrii. Zaznacz utratę danych, jeśli wystąpiła.}
\section{Pomiary na Raspberry Pi}
\writingnote{Przedstaw opóźnienia, CPU, RAM i rozmiary modeli ze wskazaniem procedury pomiaru.}
\section{Przykłady błędnych i poprawnych detekcji}
\writingnote{Pokaż wybrane przebiegi czasowe z etykietami, wynikiem detektora i momentem alertu. Uzasadnij wybór przykładów.}
